\section{Related Work}

\subsection{AI-Simulated Expert Review}

The use of LLMs to simulate peer review has attracted growing attention as a mechanism for scaling evaluation capacity. Shah and Wein \cite{shah2023shah2023} investigate whether LLMs can replicate human reviewer judgments on academic manuscripts, finding moderate agreement on surface-level criteria (clarity, presentation) but weaker alignment on deeper assessments of novelty and significance. Liu et al. \cite{liu2023reviewerbot} present ReviewerBot, a system that generates structured academic reviews by conditioning on paper content and area-specific rubrics, demonstrating that domain-targeted prompting improves the specificity of feedback relative to generic review generation. The broader ``LLM-as-judge'' paradigm \cite{zheng2023judging} has shown that LLM evaluators can achieve reasonable agreement with human panels on pairwise preference tasks, though with known biases toward verbosity and self-generated content.

Most directly relevant to the present work is \citet{dellaLibera2025profiles}, which introduces token-efficient persistent reviewer profiles for academic paper evaluation. That work treats profiles as reusable artifacts with structured schemas, demonstrating that explicit persona documentation reduces per-call token overhead while preserving reviewer consistency across sessions. However, the academic review setting shares a critical property: evaluation criteria are relatively standardized. Conference rubrics specify what ``novelty,'' ``soundness,'' and ``presentation'' mean, and the broader community has converged on shared definitions over decades of practice. This standardization makes academic review a tractable first testbed but limits how far findings generalize.

Creative domain evaluation presents a structurally different challenge. There is no shared rubric for what makes a puzzle elegant, a level design satisfying, or a narrative resonant. Quality is operationalized through expert judgment, and expert judgment is itself the object of inquiry. Domain knowledge functions not as background context for applying a known rubric, but as the mechanism through which evaluation criteria are derived in the first place. We extend the AI-simulated review paradigm to this harder setting, showing that profile specificity --- not merely profile presence --- is the lever that matters when criteria cannot be assumed.

\subsection{Persona-Based Prompting}

Persona-based prompting has a substantial prior literature across several application areas. Zhang et al. \cite{zhang2018personas} demonstrate that injecting persona descriptions into dialogue systems improves response coherence and user engagement, establishing that LLMs benefit from explicit identity framing rather than generic instruction. Park et al. \cite{park2023generative} extend this to generative agents with persistent identity representations, showing that memory-augmented personas can sustain coherent behavior across multi-turn interactions in simulated social environments. Dang et al. \cite{dang2023personas} apply persona conditioning to creative writing assistance, finding that distinct persona instantiations produce meaningfully different stylistic outputs from the same model, with implications for diversity in co-creative systems.

Across this literature, persona is treated primarily as a \emph{binary} variable: either a persona description is present or it is not. The implicit assumption is that any reasonable persona specification will activate the relevant prior knowledge in the underlying model. What remains unstudied is persona quality as an independent variable --- whether the richness, specificity, and internal coherence of a persona description materially affects the quality of the output it produces, and whether this effect can be measured and replicated.

This gap is the central question of the present work. We treat persona profiles as design artifacts with measurable quality levels, comparing a baseline condition (no profile), a generic expert condition (v1: 12 profiles, $\sim$50 lines each), and a domain-specific condition (v2: 29 profiles, $\sim$80--120 lines each, with named analytical frameworks). This design allows us to decompose the contribution of persona presence from the contribution of persona quality, and to characterize the dimension along which quality improvements manifest --- not primarily in score magnitude, but in the specificity and transferability of the evaluative frameworks introduced.

\subsection{Creative Evaluation Methodology}

Evaluating creative artifacts is a foundational problem in computational creativity research. Colton and Wiggins \cite{colton2012computational} survey the field's evaluation challenges, arguing that creativity assessment requires distinguishing novelty, quality, and intentionality --- dimensions that resist reduction to scalar metrics and that depend on the evaluator's cultural and domain context. Ritchie \cite{ritchie2007assessing} proposes a formal criterion-based framework for assessing computational creativity systems, but notes that any such framework must be instantiated with domain-specific quality standards rather than universal ones. Boden \cite{boden2004creative} distinguishes combinational, exploratory, and transformational creativity, with the latter category in particular resisting automated evaluation because the evaluator must recognize when a new conceptual space has been opened.

Within game design, quality criteria are contested but not unarticulculated. Sicart \cite{sicart2011against} argues that proceduralist design frameworks --- rubrics centered on rule systems and optimal play --- miss the ethical and aesthetic dimensions of gameplay experience. Juul \cite{juul2005half} develops a theory of games as half-real, emphasizing the interaction between formal rule structures and player-projected fiction. These perspectives suggest that game design evaluation requires sensitivity to both structural properties (is the mechanism sound?) and experiential ones (does it produce the intended affect?).

Puzzle design in particular has a practitioner tradition of quality criteria that is more explicit than most creative domains. The Riven Standard, articulated by Rand Miller in reflecting on the \emph{Myst} sequel, holds that a well-designed puzzle is one that \emph{is} what the domain does: its mechanic should be inseparable from its setting, not grafted onto it \cite{miller2016riven}. Snyder \cite{snyder2019craftsmanship} enumerates craftsmanship criteria for crossword construction that translate across puzzle types: every element should pull its weight, solutions should be uniquely determined, and the solver's ``aha'' should be earned rather than forced. Our six-dimension scoring rubric (Clarity, Solvability, Elegance, Reading Reward, Fun, Confirmation) operationalizes these practitioner principles into measurable dimensions while preserving their emphasis on the solver's experiential arc. The rubric's contribution is not to replace qualitative judgment but to give AI reviewer panels a stable coordinate system within which qualitative observations can be located and compared.

\subsection{Puzzle Hunt as a Research Domain}

Puzzle hunts occupy an unusual position in the landscape of creative artifacts: they are both genuinely creative and partially verifiable. Each puzzle has a single correct answer, so correctness can be confirmed objectively. But whether a puzzle is \emph{good} --- fair, elegant, thematically resonant, and appropriately difficult --- is a matter of expert judgment that the puzzle hunt community has developed sophisticated vocabulary for debating. This combination of verifiable correctness and contested quality makes puzzle hunt design a particularly tractable testbed for empirical work on creative evaluation.

Several properties distinguish puzzle hunts from superficially similar domains. Unlike trivia, where domain knowledge determines whether a solver can answer a question, puzzle hunt design requires that domain knowledge \emph{be the mechanism}: the puzzle should be impossible without understanding the domain, but fully solvable once one does. This structural property --- sometimes called thematic integration --- is itself a quality dimension that skilled evaluators assess and that naive evaluators tend to miss \cite{williams2014puzzlehunt}. Unlike narrative fiction or visual art, hunts have a documented expert community (Mystery Hunt constructors, BAPHL designers, MITMH veterans) whose evaluative philosophies are discussed publicly in post-hunt retrospectives, making it possible to construct profiles grounded in real critical perspectives rather than invented ones.

The empirical study of puzzle design as creative practice is small but growing \cite{link2018puzzledesign, ormond2022huntcraft}. Prior work has examined solver behavior, difficulty calibration, and hint design, but has not examined AI-assisted evaluation or the role of reviewer persona in shaping evaluation quality. Our contribution is to establish puzzle hunt design as a controlled creative evaluation testbed: one with sufficient domain structure to support empirical comparison, sufficient creative latitude to make evaluation non-trivial, and sufficient community documentation to ground AI reviewer profiles in authentic expert perspectives.
